% Figure: OUR vs OTR during a batch culture: OTR (capacity) vs OUR
% (demand), oxygen-limitation transition.
\begin{figure}[htbp]
\centering
\definecolor{engorange}{RGB}{217, 119, 6}
\begin{tikzpicture}[
  axis/.style={->, thick},
  our/.style={engred, very thick},
  otr/.style={engblue, very thick},
  dox/.style={engorange, very thick},
  shade/.style={engred!10},
  lab/.style={font=\small},
]

\draw[axis] (0, 0) -- (10, 0) node[right, lab] {time (h)};
\draw[axis] (0, 0) -- (0, 5.5) node[above, lab] {rate (mmol O$_2$ / L / h)};

% Time ticks
\foreach \x/\v in {0/0, 2/4, 4/8, 6/12, 8/16}
  {\draw (\x, 0) -- (\x, -0.08); \node[lab, anchor=north, font=\scriptsize] at (\x, -0.08) {\v};}
% y ticks
\foreach \y/\v in {0/0, 1/20, 2/40, 3/60, 4/80, 5/100}
  {\draw (0, \y) -- (-0.08, \y); \node[lab, anchor=east, font=\scriptsize] at (-0.08, \y) {\v};}

% OTR capacity (rises with agitation; here roughly flat at ~3.5)
\draw[otr] (0, 3.5) -- (9.5, 3.5);
\node[lab, engblue, anchor=west] at (9.5, 3.55) {OTR (capacity)};

% OUR (rises exponentially in lag/exponential, plateau when OTR-limited)
\draw[our] plot[smooth, samples=80, domain=0:5.5]
  (\x, {0.4 * exp(0.45 * \x)});
\draw[our] (5.5, 3.5) -- (9, 3.5);
\node[lab, engred, anchor=west] at (9, 3.6) {OUR (demand)};

% Dissolved O_2 trace on secondary scale
\draw[dox] plot[smooth, samples=80, domain=0:4.5]
  (\x, {5.0 - 0.2 * \x});
\draw[dox] plot[smooth, samples=80, domain=4.5:9]
  (\x, {5.0 - 0.2 * 4.5 - 1.6 * (\x - 4.5) / (9 - 4.5)});
\node[lab, engorange, anchor=west] at (9, 2.4) {$C_L$ (D.O.)};

% Shade oxygen-limited region
\fill[shade, opacity=0.6] (5.5, 0) rectangle (9.5, 3.5);
\node[lab, anchor=south, engred] at (7.5, 0.1) {oxygen-limited};

\end{tikzpicture}
\caption{Oxygen uptake rate (OUR, red) and oxygen transfer rate (OTR,
blue) during an aerobic batch fermentation. OUR rises with biomass
during lag and exponential growth; OTR is set by mixing, gas flow, and
$k_L a$ for the reactor as configured. When the rising OUR meets the
OTR ceiling, the dissolved oxygen $C_L$ (orange) collapses from the
saturating value and the culture becomes oxygen-limited. The
production engineer's job is to push the OTR ceiling up (more
agitation, more gas, higher pressure, oxygen enrichment) or to constrain
the OUR (lower temperature, fed-batch substrate limitation, lower
inoculum) so that the rates meet at a tolerable dissolved-oxygen
setpoint. The figure axes carry order-of-magnitude rates for a small
production-scale microbial fermenter.}
\label{fig:vol06ch10:our-otr}
\end{figure}
